\section{Attention models for price forecasting --- what they do, what they cost}

\subsection{The motivation}

LEAR and LightGBM encode the past through hand-crafted lags: prices at
lag 24, 48, 168 hours. The 168-hour lag reaches back one week. If price
depends on a regime that formed a month ago (a prolonged high-gas period,
a cold snap, a nuclear outage cluster), those lags are blind to it.
Attention models can look back months, and do so dynamically:
every forecast hour decides for itself which past hours matter most.

\subsection{Self-attention in one sentence}

Given a sequence of vectors $z_1, \ldots, z_T$, attention computes for
each position $t$ a weighted average of all other positions. The weight
between $t$ and $s$ is proportional to $\exp(q_t \cdot k_s / \sqrt{d})$,
where $q_t$ (``query'') and $k_s$ (``key'') are learned linear projections.
Positions that are semantically similar attract large weights.

Cost: $O(T^2 d)$. For $T = 2016$ (84 days hourly), this is $4 \times 10^6$
operations per layer --- feasible on GPU/MPS but slow.

\subsection{Temporal Fusion Transformer (TFT)}

TFT (Lim et al.\ 2021) wraps attention inside three structures.

\textbf{Variable Selection Network (VSN).} Before the sequence enters
the LSTM, the VSN computes a soft weight for each input feature at each
timestep. This is the key interpretability contribution: export the
weights and you know which feature mattered at which hour, without any
post-hoc explanation method. In our price model, the VSN learns to
up-weight the RES forecast hours before sunset and the TSO load forecast
on cold winter mornings.

\textbf{LSTM encoder-decoder.} The long series goes through an LSTM
first; this compresses local patterns (within-day seasonality, hour-to-hour
autocorrelation) into a fixed-size state. Only the resulting hidden states
enter the attention layer --- the sequence is shorter, so attention is cheaper.

\textbf{Temporal self-attention.} Attention over the compressed LSTM
states. This is where the long-range dependencies are captured.
Quantile output heads ($\tau = 0.1, 0.5, 0.9$) project to the three
forecast values.

\subsection{What our TFT test found}

We ran a direct test: does longer context help? Encoder lengths
1 / 4 / 12 weeks, otherwise identical. Results on the price task
(same 4320-hour window):

\begin{center}
\begin{tabular}{lcc}
model & rMAE & spike MAE \\
\hline
LEAR + conformal & 0.66 & 52.2 \\
LightGBM + conformal & 0.67 & 49.8 \\
TFT, 12-week context & 0.86 & 57.8 \\
TFT, 4-week context & 0.88 & 54.3 \\
TFT, 1-week context & 0.93 & 72.4 \\
\end{tabular}
\end{center}

Two things are both true:

\begin{itemize}
  \item Long context works: rMAE falls monotonically as the encoder
        looks further back. The attention mechanism IS extracting
        something from month-old history.
  \item It is not enough: the best TFT trails LEAR by 30\%. The binding
        limit is not architecture but information --- the outage and
        fuel features that explain spikes are not in our input set.
        A perfect attention model over bad inputs cannot beat a linear
        model over the right inputs.
\end{itemize}

\subsection{PatchTST --- cheaper long context via patching}

PatchTST (Nie et al.\ 2023) cuts the input series into overlapping
\emph{patches} (e.g.\ 24-hour windows, 12-hour stride). Each patch
becomes one transformer token. For a 2016-hour encoder with 24-hour
patches and 12-hour stride, the sequence length is
$(2016 - 24) / 12 + 1 = 166$ tokens instead of 2016 --- 12$\times$ shorter.
Attention cost drops from $2016^2$ to $166^2$ (144$\times$ faster).

The price paid: within-patch structure is lost. If two adjacent hours
differ in ways the patch mean masks (a sudden spike, a DST gap),
the model is blind to that.

Channel independence: each input variable passes through its own
copy of the same weights, with no cross-channel attention. Our data
supports this (1 of 21 pairs $|\rho| > 0.5$). The bet is that
price channels are weakly coupled at the hourly level, so sharing
weights across channels is better than trying to learn sparse interactions
from a small dataset.

\subsection{VSN vs SHAP vs ablation --- three interpretation layers}

Section \ref{sec:shap-vs-ablation} compared SHAP and ablation. For TFT,
a third signal exists: the VSN weights.

\begin{center}
\begin{tabular}{lll}
Method & What it measures & Cost \\
\hline
VSN weight & in-model softmax over features at each hour & free (a forward pass) \\
SHAP & local gradient of output w.r.t.\ input & 1 min per day \\
Ablation & marginal skill contribution of feature groups & one backtest per group \\
\end{tabular}
\end{center}

They agree on the direction but differ in granularity. VSN is the cheapest
and the only one that varies hour-by-hour. SHAP is the most trusted for
tabular models. Ablation is the most honest for groups.

Use all three when a model is in production; use VSN alone for daily
driver reports (it is free).

\subsection{Interview line}

``We tested the long-context hypothesis directly: 12 weeks of attention
beats 1 week --- the mechanism is real --- but LEAR with the right
features still wins by 30\%. We concluded the bottleneck was missing
information (outages, fuel), not architecture. So we added the fuel
features first, then ran TFT again with HPO. The honest comparison is
in the model card.''
